\section{Lecture 5: Inference and Prediction}

\subsection{Guidelines for Transformations}

Finding the appropriate transformation depends on the problem, and there is no universal solution. However, several general guidelines are useful when diagnosing linear regression models.

If diagnostic plots suggest violations of the linearity assumption, transforming the predictors may help achieve a more linear relationship between the response and predictors.

If residual plots indicate non-constant variance or departures from normality, transforming the response variable can help stabilize the variance and improve the distribution of the residuals.

\subsection{Residual Plots with Categorical Predictors}

When residuals are plotted against categorical predictors, the resulting visualization consists of boxplots. The same diagnostic principles apply: residuals should be centered around zero, exhibit similar variability across groups, and show no systematic patterns.

\subsubsection{Example: Residual Diagnostics for the Iris Data}

We consider the linear model
\[
\text{Petal.Length} \sim \text{Species} + \text{Sepal.Length}.
\]

\begin{figure}[H]
    \centering
    \includegraphics[width=0.8\linewidth]{pictures/lec5.png}
    \caption{Linearity satisfied, Normality satisfied, Constant variance not satisfied.}
\end{figure}

\clearpage

\subsection{Review: Sampling Distributions for Simple Linear Regression}

For simple linear regression, the sampling distributions of the least squares estimators are
\[
\bhat{\beta}_0 \mid X
\sim
\mathcal{N}
\!\left(
\beta_0,\,
\sigma^2
\left(
\frac{1}{n}
+
\frac{\bar{x}^2}{S_{xx}}
\right)
\right),
\qquad
\bhat{\beta}_1 \mid X
\sim
\mathcal{N}
\!\left(
\beta_1,\,
\frac{\sigma^2}{S_{xx}}
\right),
\]
where
\[
S_{xx} = \sum_{i=1}^n (x_i - \bar{x})^2.
\]

Standardizing these estimators yields
\[
Z
=
\frac{\bhat{\beta}_0 - \beta_0}
{\sigma \sqrt{\frac{1}{n} + \frac{\bar{x}^2}{S_{xx}}}}
\sim \mathcal{N}(0,1),
\qquad
Z
=
\frac{\bhat{\beta}_1 - \beta_1}
{\sigma / \sqrt{S_{xx}}}
\sim \mathcal{N}(0,1).
\]

For simple linear regression ($p = 1$), the vector of estimators satisfies
\[
\bbeta
=
\begin{pmatrix}
\bhat{\beta}_0 \\
\bhat{\beta}_1
\end{pmatrix}
\sim
\mathcal{N}_2
\!\left(
\bbeta,\,
\sigma^2 (X^\top X)^{-1}
\right).
\]

If $\sigma$ were known, these standardized statistics could be used directly to perform hypothesis tests and construct confidence intervals for $\beta_0$ and $\beta_1$. In practice, however, the standard deviation $\sigma$ is usually unknown.

The appropriate remedy is to replace $\sigma^2$ with its unbiased estimator and use the t-test:
\[
S^2
=
\hat{\sigma}^2
=
\frac{1}{n - p - 1}
\sum_{i=1}^n \hat{e}_i^2
=
\frac{\RSS}{n - p - 1}.
\]

For simple linear regression ($p = 1$), this simplifies to
\[
S^2
=
\hat{\sigma}^2
=
\frac{1}{n - 2}
\sum_{i=1}^n
\left(
y_i - \bhat{\beta}_0 - \bhat{\beta}_1 x_i
\right)^2.
\]

\subsection{Inferences About Regression Parameters}

\subsubsection{A) Hypothesis Testing Framework}

Let $\theta$ denote a generic regression parameter (e.g., $\beta_0$ or $\beta_1$). A hypothesis test is defined by
\[
H_0 : \theta = \theta_0
\qquad \text{versus} \qquad
H_1 : \theta \neq \theta_0,
\]
where $\theta_0$ is a specified value under the null hypothesis.

The test is conducted using an appropriate test statistic, whose distribution under $H_0$ is known.

\subsubsection{B) Confidence Interval Framework}

Let $\theta$ denote a parameter of interest. A confidence interval for $\theta$ is given by
\[
\hat{\theta}
\;\pm\;
(\text{critical value}) \times \se(\hat{\theta}),
\]
where the critical value is determined by the appropriate sampling distribution (t or Z test).

\subsection{Test Statistic, p-Value, and Rejection Rule for SLR}

To test
\[
H_0 : \beta_1 = \beta_1^\star
\qquad \text{versus} \qquad
H_1 : \beta_1 \neq \beta_1^\star,
\]
the test statistic is
\[
T
=
\frac{\bhat{\beta}_1 - \beta_1}
{\hat{\sigma} / \sqrt{S_{xx}}}
=
\frac{\bhat{\beta}_1 - \beta_1}
{\se(\bhat{\beta}_1)} \sim t_{n-2},
\]
and the observed test statistic is
\[
T_{\text{obs}}
=
\frac{\bhat{\beta}_1 - \beta_1^\star}
{\se(\bhat{\beta}_1)},
\]
where $\se(\bhat{\beta}_1)$ is the estimated standard error of $\bhat{\beta}_1$ (as reported by \texttt{R}).

\vspace{0.5em}
\hrule
\vspace{0.5em}

To test
\[
H_0 : \beta_0 = \beta_0^\star
\qquad \text{versus} \qquad
H_1 : \beta_0 \neq \beta_0^\star,
\]
the test statistic is
\[
T
=
\frac{\bhat{\beta}_0 - \beta_0}
{\hat{\sigma} \sqrt{\frac{1}{n} + \frac{\bar{x}^2}{S_{xx}}}}
=
\frac{\bhat{\beta}_0 - \beta_0}
{\se(\bhat{\beta}_0)} \sim t_{n-2},
\]
and the observed test statistic is
\[
T_{\text{obs}}
=
\frac{\bhat{\beta}_0 - \beta_0^\star}
{\se(\bhat{\beta}_0)},
\]
where $\se(\bhat{\beta}_0)$ is the estimated standard error of $\bhat{\beta}_0$ (as reported by \texttt{R}).

\vspace{0.5em}
\hrule
\vspace{0.5em}

\subsubsection{Test Hypothesis}

The two-sided p-value is
\[
\text{p-value}
=
2 \, \mathbb{P}
\bigl(
|t_{n-2}| \ge |T_{\text{obs}}|
\bigr).
\]

The test rejects $H_0$ at significance level $\alpha$ if
\[
\text{p-value} < \alpha,
\quad \text{OR} \quad
|T_{\text{obs}}| \ge t_{\alpha/2,\,n-2}.
\]

Here, $t_{\alpha/2,\,n-2}$ is the $100(1-\alpha/2)$th percentile of the $t_{n-2}$ distribution, i.e.,
\[
\mathbb{P}\bigl(t_{n-2} > t_{\alpha/2,\,n-2}\bigr) = \alpha/2.
\]

\subsubsection{Confidence Interval}

A $100(1-\alpha)\%$ confidence interval for $\beta_1$ is given by
\[
\bigl(
\bhat{\beta}_1 - t_{\alpha/2,\,n-2}\,\se(\bhat{\beta}_1),
\;
\bhat{\beta}_1 + t_{\alpha/2,\,n-2}\,\se(\bhat{\beta}_1)
\bigr).
\]

Here, $t_{\alpha/2,\,n-2}$ is the $100(1-\alpha/2)$th quantile of the $t$-distribution with $n-2$ degrees of freedom.

\vspace{0.5em}
\hrule
\vspace{0.5em}

A $100(1-\alpha)\%$ confidence interval for $\beta_0$ is given by
\[
\bigl(
\bhat{\beta}_0 - t_{\alpha/2,\,n-2}\,\se(\bhat{\beta}_0),
\;
\bhat{\beta}_0 + t_{\alpha/2,\,n-2}\,\se(\bhat{\beta}_0)
\bigr).
\]

Here, $t_{\alpha/2,\,n-2}$ is the $100(1-\alpha/2)$th quantile of the $t$-distribution with $n-2$ degrees of freedom.

\subsection{Clarifications on p-Value and Test Hypothesis}

\begin{itemize}
\item \textbf{Definition 1:} It is the smallest level of significance at which $H_0$ would be rejected based on the observed data.
\item \textbf{Definition 2:} It is the probability of observing the result as or more extreme than that actually observed if $H_0$ is true.
\item In naive terms, the p-value suggests how surprising the observed sample is if we assume $H_0$ to be true.
\item Conventionally, we compare the p-value to 0.01, 0.05, or 0.1. If the p-value is less than a predefined cut-off, we reject $H_0$.
\end{itemize}

The hypothesis test rejects $H_0$ at significance level $\alpha$ if
\[
\text{p-value} < \alpha,
\]
which is equivalent to
\[
|T_{\text{obs}}| \ge t_{\alpha/2,\,n-2}.
\]

\subsubsection{Visualization and Meaning of $t_{0.1,10}$}

\[
t_{0.10,\,10}
\;\overset{\text{R}}{=}\;
\texttt{qt}(0.90, \text{df} = 10).
\]

\begin{figure}[H]
\centering
\includegraphics[width=0.5\textwidth]{pictures/t_distribution_df10.png}
\end{figure}

\subsection{Approach 1: Critical Region Approach ($\alpha$ Known Beforehand)}

\textbf{Idea.} Before observing the data, a significance level $\alpha$ is fixed. This choice determines the rejection regions in the tails of the $t_{n-2}$ distribution.

\textbf{Decision rule.} Compute the test statistic $T_{\text{obs}}$. Reject $H_0$ if
\[
|T_{\text{obs}}| \ge t_{\alpha/2,\,n-2}.
\]

\textbf{Interpretation.} The shaded tail regions have total probability $\alpha$ under $H_0$. If the observed test statistic falls into these regions, it is considered too extreme to be consistent with the null hypothesis.

\begin{figure}[H]
    \centering
    \includegraphics[width=0.5\linewidth]{pictures/p_value.png}
\end{figure}

\subsection{Approach 2: p-Value Approach ($\alpha$ Not Fixed Beforehand)}

\textbf{Idea.} First compute how extreme the observed test statistic is, assuming $H_0$ is true. This extremeness is summarized by the p-value.

\textbf{Definition.} For a two-sided test,
\[
\text{p-value}
=
2\,\mathbb{P}\bigl(|t_{n-2}| \ge |T_{\text{obs}}|\bigr).
\]

\textbf{Decision rule.} After computing the p-value, reject $H_0$ for any chosen significance level $\alpha$ if
\[
\text{p-value} < \alpha.
\]

\textbf{Interpretation.} The p-value measures the strength of evidence against $H_0$. Smaller p-values indicate stronger evidence against the null hypothesis.

\begin{figure}[H]
    \centering
    \includegraphics[width=0.5\linewidth]{pictures/critical_region.png}
\end{figure}

\clearpage

\subsection{SLR Example: Finding CI and Testing Hypothesis Using R Output}

\begin{figure}[H]
    \centering
    \includegraphics[width=\linewidth]{pictures/lec5outputexample.png}
\end{figure}

The 95\% confidence interval for $\beta_0$ is
\[
\bhat{\beta}_0 \pm t_{0.975,\,1076}\,\se(\bhat{\beta}_0)
=
33.8866 \pm 1.96(1.8324)
=
(30.30,\; 37.47).
\]

The 95\% confidence interval for $\beta_1$ is
\[
\bhat{\beta}_1 \pm t_{0.975,\,1076}\,\se(\bhat{\beta}_1)
=
0.5141 \pm 1.96(0.0271)
=
(0.461,\; 0.567).
\]

To test $H_0:\beta_0=0$, the observed test statistic is
\[
T_{\text{obs}} = \frac{33.8866}{1.8324} = 18.49.
\]
Since the p-value is $<2\times10^{-16}$ and $|T_{\text{obs}}| > t_{0.025,1076}$, we reject $H_0$.

To test $H_0:\beta_1=0$, the observed test statistic is
\[
T_{\text{obs}} = \frac{0.5141}{0.0271} = 19.01.
\]
Since the p-value is $<2\times10^{-16}$ and $|T_{\text{obs}}| > t_{0.025,1076}$, we reject $H_0$.

These confidence intervals agree with the output produced directly in \textsf{R}:

\begin{figure}[H]
    \centering
    \includegraphics[width=0.5\linewidth]{pictures/lec5coef.png}
\end{figure}

\clearpage

\subsection{Confidence Intervals: Multiple Linear Regression}

For $j = 0,1,\ldots,p$,
\[
T_j
=
\frac{\bhat{\beta}_j - \beta_j}{\hat{\sigma}_{\bhat{\beta}_j}}
\;\sim\;
t_{n-p-1},
\]
where
\[
\hat{\sigma}_{\bhat{\beta}_j}
=
\hat{\sigma}
\sqrt{(\bX^\top \bX)^{-1}_{\,j+1,\,j+1}}.
\]

A $100(1-\alpha)\%$ confidence interval for $\beta_j$, $j=0,1,\ldots,p$, is given by
\[
\left(
\bhat{\beta}_j
-
t_{\alpha/2,\,n-p-1}\,\hat{\sigma}_{\bhat{\beta}_j},
\;
\bhat{\beta}_j
+
t_{\alpha/2,\,n-p-1}\,\hat{\sigma}_{\bhat{\beta}_j}
\right),
\]
where $t_{\alpha/2,\,n-p-1}$ is the $100(1-\alpha/2)$th quantile of the $t$-distribution with $n-p-1$ degrees of freedom.

\subsection{Prediction}

After fitting the regression model, simple or multiple, we observe a new sample
\[
\mathbf{x}^\star = (1, x_1^\star, \ldots, x_p^\star)^\top
\]
without observing its response.

If the new sample follows the same linear regression model, its response is
\[
Y^\star = \beta_0 + \beta_1 x_1^\star + \cdots + \beta_p x_p^\star + \varepsilon^\star,
\qquad
\varepsilon^\star \sim N(0,\sigma^2).
\]

The predicted response is
\[
\hat{Y}^\star = \mathbf{x}^{\star\top}\bbeta.
\]

\subsubsection{Distribution of $\hat{Y}^\star$}

\[
\bbeta \sim N_{p+1}\!\left(
\bbeta,\,
\sigma^2(\bX^\top\bX)^{-1}
\right).
\]

The predicted value $\hat{Y}^\star$ is normally distributed with
\[
\E(\hat{Y}^\star)
=
\mathbf{x}^{\star\top}\E(\bbeta)
=
\mathbf{x}^{\star\top}\bbeta,
\]
and
\[
\Var(\hat{Y}^\star)
=
\Var(\mathbf{x}^{\star\top}\bbeta)
=
\mathbf{x}^{\star\top}
\Var(\bbeta)
\mathbf{x}^\star
=
\sigma^2\,\mathbf{x}^{\star\top}(\bX^\top\bX)^{-1}\mathbf{x}^\star.
\]

Hence,
\[
\hat{Y}^\star
\sim
N\!\left(
\mathbf{x}^{\star\top}\bbeta,\,
\sigma^2\,\mathbf{x}^{\star\top}(\bX^\top\bX)^{-1}\mathbf{x}^\star
\right).
\]

\clearpage

\subsubsection{SLR: Distribution of $\hat{Y}^\star$ and Inference}

In simple linear regression,
\[
\bbeta
=
\begin{pmatrix}
\beta_0\\
\beta_1
\end{pmatrix},
\qquad
\mathbf{x}^\star
=
\begin{pmatrix}
1\\
x^\star
\end{pmatrix}.
\]

Then
\[
\mathbf{x}^{\star\top}\bbeta
=
\beta_0 + \beta_1 x^\star,
\]
and
\[
\mathbf{x}^{\star\top}(\bX^\top\bX)^{-1}\mathbf{x}^\star
=
\frac{1}{n}
+
\frac{(\bar{x}-x^\star)^2}{S_{xx}}.
\]

Thus,
\[
\hat{Y}^\star
\sim
N\!\left(
\beta_0 + \beta_1 x^\star,\,
\sigma^2\!\left(
\frac{1}{n}
+
\frac{(\bar{x}-x^\star)^2}{S_{xx}}
\right)
\right).
\]

The estimated standard error of $\hat{Y}^\star$ is
\[
\widehat{\se}(\hat{Y}^\star)
=
\hat{\sigma}
\sqrt{
\frac{1}{n}
+
\frac{(\bar{x}-x^\star)^2}{S_{xx}}
}.
\]

Standardization:
\[
\frac{
\hat{Y}^\star - (\beta_0 + \beta_1 x^\star)
}{
\sigma
\sqrt{
\frac{1}{n}
+
\frac{(\bar{x}-x^\star)^2}{S_{xx}}
}
}
\;\sim\;
N(0,1).
\]

Replacing $\sigma$ with $\hat{\sigma}$:
\[
\frac{
\hat{Y}^\star - (\beta_0 + \beta_1 x^\star)
}{
\hat{\sigma}
\sqrt{
\frac{1}{n}
+
\frac{(\bar{x}-x^\star)^2}{S_{xx}}
}
}
\;\sim\;
t_{n-2}.
\]

\begin{keyresult}[Confidence Interval for the Mean Response (SLR)]
A $100(1-\alpha)\%$ confidence interval for $\E(Y^\star \mid x^\star) = \beta_0 + \beta_1 x^\star$ is
\[
\hat{Y}^\star
\pm
t_{\alpha/2,\,n-2}\,
\hat{\sigma}
\sqrt{
\frac{1}{n}
+
\frac{(\bar{x}-x^\star)^2}{S_{xx}}
}.
\]
\end{keyresult}

\subsubsection{MLR: Distribution of $\hat{Y}^\star$ and Inference}

\[
\hat{Y}^\star
\sim
N\!\left(
\mathbf{x}^{\star\top}\bbeta,
\;
\mathbf{x}^{\star\top}\sigma^2(\bX^\top\bX)^{-1}\mathbf{x}^\star
\right).
\]

\textbf{Standardize:}
\[
\frac{
\hat{Y}^\star - \mathbf{x}^{\star\top}\bbeta
}{
\sigma
\sqrt{
\mathbf{x}^{\star\top}(\bX^\top\bX)^{-1}\mathbf{x}^\star
}
}
\;\sim\;
N(0,1).
\]

\textbf{Replacing $\sigma$ with $\hat{\sigma}$:}
\[
\frac{
\hat{Y}^\star - \mathbf{x}^{\star\top}\bbeta
}{
\hat{\sigma}
\sqrt{
\mathbf{x}^{\star\top}(\bX^\top\bX)^{-1}\mathbf{x}^\star
}
}
\;\sim\;
t_{n-p-1}.
\]

\begin{keyresult}[Confidence Interval for the Mean Response (MLR)]
A $100(1-\alpha)\%$ confidence interval for the mean response of $Y^\star$ (given $\mathbf{x}^\star$) is
\[
\hat{Y}^\star
\pm
t_{\alpha/2,\,n-p-1}\,
\hat{\sigma}
\sqrt{
\mathbf{x}^{\star\top}(\bX^\top\bX)^{-1}\mathbf{x}^\star
}.
\]
\end{keyresult}

\subsection{Prediction Interval for the Actual Value $Y^\star$ (SLR)}

Assume
\[
\varepsilon^\star \sim \mathcal{N}(0,\sigma^2)
\]
and $\varepsilon^\star$ is independent of $\varepsilon_1,\dots,\varepsilon_n$.

The prediction error is
\[
Y^\star - \hat{Y}^\star.
\]

Its expectation is
\[
\E(Y^\star - \hat{Y}^\star) = 0,
\]
and its variance is
\[
\Var(Y^\star - \hat{Y}^\star)
= \Var(Y^\star) +\Var(\hat{Y}^\star) =
\sigma^2
\left(
1 + \frac{1}{n} + \frac{(\bar{x}-x^\star)^2}{S_{xx}}
\right).
\]

Thus,
\[
Y^\star - \hat{Y}^\star
\sim
\mathcal{N}\!\left(
0,
\;
\sigma^2
\left(
1 + \frac{1}{n} + \frac{(\bar{x}-x^\star)^2}{S_{xx}}
\right)
\right).
\]

Standardizing,
\[
\frac{
Y^\star - \hat{Y}^\star
}{
\sigma
\sqrt{
1 + \frac{1}{n} + \frac{(\bar{x}-x^\star)^2}{S_{xx}}
}
}
\sim
\mathcal{N}(0,1).
\]

Replacing $\sigma$ with $\hat{\sigma}$,
\[
\frac{
Y^\star - \hat{Y}^\star
}{
\hat{\sigma}
\sqrt{
1 + \frac{1}{n} + \frac{(\bar{x}-x^\star)^2}{S_{xx}}
}
}
\sim
t_{n-2}.
\]

\begin{keyresult}[Prediction Interval for the Actual Value (SLR)]
A $100(1-\alpha)\%$ prediction interval for the actual value $Y^\star$ is
\[
\hat{Y}^\star
\pm
t_{\alpha/2,\,n-2}
\;
\hat{\sigma}
\sqrt{
1 + \frac{1}{n} + \frac{(\bar{x}-x^\star)^2}{S_{xx}}
}.
\]
\end{keyresult}

\subsection{Prediction Interval for the Actual Value $Y^\star$ (MLR)}

The prediction error satisfies
\[
Y^\star - \hat{Y}^\star
\sim
\mathcal{N}\!\left(
0,
\;
\sigma^2
\left(
1 +
\mathbf{x}^{\star\top}
(\bX^\top \bX)^{-1}
\mathbf{x}^\star
\right)
\right).
\]

The estimated standard error is
\[
\widehat{\se}(Y^\star - \hat{Y}^\star)
=
\hat{\sigma}
\sqrt{
1 +
\mathbf{x}^{\star\top}
(\bX^\top \bX)^{-1}
\mathbf{x}^\star
}.
\]

\begin{keyresult}[Prediction Interval for the Actual Value (MLR)]
A $100(1-\alpha)\%$ prediction interval for $Y^\star$ is
\[
\hat{Y}^\star
\pm
t_{\alpha/2,\,n-p-1}
\;
\hat{\sigma}
\sqrt{
1 +
\mathbf{x}^{\star\top}
(\bX^\top \bX)^{-1}
\mathbf{x}^\star
}.
\]
\end{keyresult}

\clearpage

\subsection{Summary: Confidence vs Prediction Intervals (SLR)}

For a new observation,
\[
Y^\star = \beta_0 + \beta_1 x^\star + \varepsilon^\star,
\qquad
\varepsilon^\star \sim \mathcal{N}(0,\sigma^2).
\]

\subsubsection{95\% Confidence Interval (CI) for the Mean Response}

This interval estimates the mean response
\[
\E(Y^\star \mid x^\star) = \beta_0 + \beta_1 x^\star.
\]

\[
\hat{Y}^\star
\;\pm\;
t_{0.025,\,n-2}\;
\hat{\sigma}
\sqrt{
\frac{1}{n}
+
\frac{(\bar{x} - x^\star)^2}{S_{xx}}
}.
\]

\subsubsection{95\% Prediction Interval (PI) for the Actual Value $Y^\star$}

This interval predicts the actual random outcome $Y^\star$.

\[
\hat{Y}^\star
\;\pm\;
t_{0.025,\,n-2}\;
\hat{\sigma}
\sqrt{
\boxed{1}
+
\frac{1}{n}
+
\frac{(\bar{x} - x^\star)^2}{S_{xx}}
}.
\]

\clearpage
